\chapter{Radiology}

\medterm{Computed Tomography (CT)}-- Cross-sectional imaging using rotating X-ray tube and detectors. Produces axial images, reconstructable into sagittal, coronal, 3D. CT number measured in Hounsfield units (HU): water = 0, air = -1000, bone = +1000. Contrast: iodine-based, enhances vascular structures and perfused organs. Speed: seconds. Radiation dose: higher than plain X-ray.

\medterm{Magnetic Resonance Imaging (MRI)}-- Uses strong magnetic field and radiofrequency pulses to generate images from hydrogen protons. Excellent soft tissue contrast. No ionizing radiation. Sequences: T1-weighted (fat bright, fluid dark), T2-weighted (fluid bright), FLAIR (suppresses CSF), DWI (diffusion restriction). Contrast: gadolinium-based. Contraindications: pacemaker, certain implants, claustrophobia.

\medterm{Ultrasound}-- Real-time imaging using high-frequency sound waves. No ionizing radiation. Useful for abdominal organs, obstetrics, vascular studies, musculoskeletal. Doppler assesses blood flow. Limited by bone and air.

\medterm{X-ray (Radiography)}-- Projectional imaging using ionizing radiation. Dense structures (bone, contrast) appear white; air appears black. Fast, inexpensive, widely available. Two-dimensional projection of three-dimensional structures. Mammography: specialized breast X-ray for screening.

\medterm{Fluoroscopy}-- Real-time X-ray imaging. Used for barium studies, angiography, catheter placement, joint injections. Continuous or pulsed radiation. Radiation dose higher than plain X-ray.

\medterm{Positron Emission Tomography (PET)}-- Functional imaging using positron-emitting radiotracers. FDG-PET most common. Metabolic activity visualized. Combined with CT (PET/CT) for anatomical correlation. Used in oncology, neurology, cardiology.

\medterm{Single Photon Emission Computed Tomography (SPECT)}-- Nuclear medicine imaging using gamma-emitting radiotracers (99mTc most common). Rotating gamma camera acquires images. Less resolution than PET but more widely available and cheaper. Used for cardiac perfusion, bone scans, brain imaging.

\medterm{Contrast Agent}-- Substance enhancing image visibility. X-ray/CT: iodine-based (positive contrast). MRI: gadolinium-based (shortens T1, bright on T1). Ultrasound: microbubble agents. Barium sulfate for GI studies.

\medterm{Interventional Radiology}-- Image-guided minimally invasive procedures. Examples: angioplasty, stenting, embolization, biopsy, ablation, drain placement. Reduces need for open surgery.

\medterm{Mammography}-- Breast X-ray for cancer screening. Digital mammography standard. Tomosynthesis (3D mammography) improves detection. BI-RADS classification system. Screening: annual or biennial for women 40-74.

\medterm{DEXA (Dual-Energy X-ray Absorptiometry)}-- Measures bone mineral density (BMD) at spine and hip. T-score compares to young adult reference. Osteoporosis: T-score <= -2.5. Osteopenia: -1.0 to -2.5. Low radiation dose.

\medterm{Angiography}-- X-ray imaging of blood vessels after contrast injection. Digital subtraction angiography (DSA) gold standard. Used for diagnosing and treating vascular diseases.

\medterm{Doppler Ultrasound}-- Ultrasound technique measuring blood flow velocity and direction. Color Doppler, spectral Doppler, power Doppler. Used for vascular studies, cardiac evaluation, organ perfusion.

\medterm{CT Angiography (CTA)}-- CT imaging of blood vessels after IV contrast. Fast, non-invasive. Used for pulmonary embolism, aortic dissection, coronary artery disease, peripheral arterial disease.

\medterm{MR Angiography (MRA)}-- Vascular imaging using MRI without ionizing radiation. Time-of-flight (TOF) or contrast-enhanced (CE-MRA). Used for cerebral aneurysms, renal artery stenosis, peripheral arterial disease.

\medterm{MRI Brain}-- Excellent for brain pathology: tumors, stroke, demyelination, infection, degeneration. Sequences: T1, T2, FLAIR, DWI, SWI, contrast-enhanced. DWI detects acute ischemia within minutes.

\medterm{MRI Spine}-- Imaging of spinal cord, nerve roots, discs, vertebrae. Excellent soft tissue contrast. Used for disc herniation, spinal stenosis, tumors, infection.

\medterm{MRI Joint}-- Evaluation of musculoskeletal structures: cartilage, ligaments, tendons, menisci, labrum. Knee: ACL, meniscus tears. Shoulder: rotator cuff, labrum. Hip: labral tears, FAI.

\medterm{Abdominal CT}-- First-line imaging for acute abdominal pain. Evaluates appendicitis, diverticulitis, bowel obstruction, trauma, renal stones, liver masses. Oral and IV contrast often used.

\medterm{Chest X-ray}-- Initial imaging for respiratory symptoms. PA and lateral views. Evaluates pneumonia, pneumothorax, pleural effusion, heart size, mediastinal masses. Limited sensitivity for early disease.

\medterm{CT Chest}-- Detailed lung and mediastinal imaging. High-resolution CT (HRCT) for interstitial lung disease. CT pulmonary angiography (CTPA) for pulmonary embolism. Lung cancer screening with low-dose CT.

\medterm{Contrast-Induced Nephropathy}-- Acute kidney injury after iodinated contrast administration. Risk factors: pre-existing renal impairment, diabetes, dehydration. Prevention: hydration, lowest contrast volume, N-acetylcysteine controversial.

\medterm{Gadolinium-Induced Nephrogenic Systemic Fibrosis}-- Rare but serious condition in patients with severe renal impairment receiving gadolinium-based MRI contrast. Causes skin thickening, organ fibrosis. Risk reduced with macrocyclic agents.

\medterm{Radiation Safety}-- ALARA principle: As Low As Reasonably Achievable. Time, distance, shielding. Pregnancy screening before imaging. Pediatric considerations: lower doses, alternative imaging when possible.

\medterm{BI-RADS}-- Breast Imaging-Reporting and Data System. Standardizes mammography reporting. Categories: 0 (incomplete), 1 (negative), 2 (benign), 3 (probably benign), 4 (suspicious), 5 (highly suggestive of malignancy), 6 (known biopsy-proven malignancy).

\medterm{LI-RADS}-- Liver Imaging-Reporting and Data System. Standardizes reporting of liver imaging in patients with cirrhosis. Used for hepatocellular carcinoma diagnosis.

\medterm{PI-RADS}-- Prostate Imaging-Reporting and Data System. Standardizes multiparametric MRI reporting for prostate cancer detection and staging.

\medterm{Interventional Procedures}-- Image-guided treatments: angioplasty, stenting, embolization, ablation, biopsy, drain placement. Advantages: minimally invasive, faster recovery, lower complication rates.

\medterm{Radiology Report}-- Document describing imaging findings and interpretation. Structure: indication, technique, findings, impression. Clear, concise language essential. Critical findings require immediate communication.

\medterm{Artifacts}-- Image distortions not representing true anatomy. Types: motion, beam hardening, metal, aliasing, partial volume. Recognition important for accurate interpretation.

\medterm{Windowing}-- Adjusting CT image brightness and contrast to optimize visualization of specific tissues. Bone window: wide window width, high level. Soft tissue window: narrow width, mid level. Lung window: wide width, low level.

\medterm{Hounsfield Units (HU)}-- CT number scale referencing water (0 HU) and air (-1000 HU). Fat: -50 to -100 HU. Water: 0 HU. Soft tissue: +20 to +50 HU. Bone: +400 to +3000 HU. Blood: +40 to +60 HU (acute), +20 to +40 HU (chronic).

\medterm{MRI Sequences}-- T1-weighted: anatomy, fat bright. T2-weighted: pathology, fluid bright. FLAIR: suppresses CSF, lesions bright. DWI: restricted diffusion bright (acute stroke, abscess). STIR: fat suppression, fluid bright. Post-contrast T1: enhancement bright.

\medterm{PET Tracers}-- 18F-FDG: glucose metabolism (cancer, infection, inflammation). 18F-fluoride: bone turnover. 11C-metionine: protein synthesis. 68Ga-DOTATATE: somatostatin receptors. 68Ga-PSMA: prostate cancer.

\medterm{Interventional Radiology}-- See medical instruments.

\medterm{Digital Subtraction Angiography (DSA)}-- Technique subtracting pre-contrast images from post-contrast images to visualize blood vessels. Gold standard for vascular imaging. Used for diagnosing and treating vascular diseases.

\medterm{Magnetic Resonance Cholangiopancreatography (MRCP)}-- MRI technique visualizing biliary and pancreatic ducts without contrast. Non-invasive alternative to ERCP. Used for choledocholithiasis, biliary strictures, pancreatic pathology.

\medterm{MR Enterography}-- MRI technique evaluating small bowel. Uses oral contrast to distend bowel. Used for Crohn disease, small bowel tumors, bleeding.

\medterm{CT Colonography}-- CT imaging of colon after bowel preparation and insufflation. Detects polyps and masses. Alternative to colonoscopy for screening. False positives common.

\medterm{Echocardiography}-- Ultrasound of heart. Transthoracic (TTE) and transesophageal (TEE) approaches. Evaluates cardiac structure, function, valves, pericardium. Doppler assesses blood flow.

\medterm{Stress Echocardiography}-- Echocardiography during exercise or pharmacological stress. Detects coronary artery disease by identifying wall motion abnormalities.

\medterm{Ventriculography}-- Imaging ventricular function. Radionuclide ventriculography (MUGA scan) accurately measures ejection fraction. Cardiac MRI also accurate.

\medterm{Bone Scan}-- Nuclear medicine imaging using Tc-99m MDP. Detects bone metastases, fractures, infection. Sensitive but not specific. Whole body imaging.

\medterm{Thyroid Scan}-- Nuclear medicine imaging using I-123 or Tc-99m. Evaluates thyroid function and nodules. Hot nodules rarely malignant. Cold nodules require further evaluation.

\medterm{Renal Scan}-- Nuclear medicine imaging using Tc-99m compounds. Evaluates renal function, obstruction, blood flow. MAG3 for function, DTPA for GFR.

\medterm{Ventilation-Perfusion (V/Q) Scan}-- Lung imaging assessing ventilation and perfusion. Mismatch indicates pulmonary embolism. Alternative to CTPA in renal impairment or contrast allergy.

\medterm{Sialography}-- Contrast imaging of salivary ducts. Evaluates obstruction, stricture, tumors. Rarely performed now.

\medterm{Myelography}-- Contrast injected into subarachnoid space for spinal imaging. Used when MRI contraindicated. CT myelography more common now.

\medterm{Arthrography}-- Contrast injected into joint for imaging. Used for shoulder, knee, hip evaluation. MR arthrography more common now.

\medterm{Hysterosalpingography}-- Contrast imaging of uterus and fallopian tubes. Evaluates infertility, uterine anomalies, tubal patency.

\medterm{Salpingography}-- Contrast imaging of fallopian tubes. Part of hysterosalpingography.

\medterm{Cholangiography}-- Contrast imaging of bile ducts. Intraoperative, percutaneous, or ERCP approach. Evaluates biliary obstruction, stones, strictures.

\medterm{Pancreatography}-- Contrast imaging of pancreatic duct. ERCP approach. Evaluates pancreatic duct abnormalities.

\medterm{Urethrography}-- Contrast imaging of urethra. Evaluates urethral strictures, trauma, anomalies.

\medterm{Cystography}-- Contrast imaging of bladder. Evaluates trauma, fistulas, vesicoureteral reflux.

\medterm{Voiding Cystourethrogram (VCUG)}-- Fluoroscopic study of bladder and urethra during voiding. Evaluates vesicoureteral reflux, posterior urethral valves.

\medterm{Phlebography}-- Contrast imaging of veins. Evaluates deep vein thrombosis, venous obstruction. Largely replaced by Doppler ultrasound.

\medterm{Lymphangiography}-- Contrast imaging of lymphatic vessels. Evaluates lymphatic obstruction, fistulas. Largely replaced by MRI/CT.

\medterm{Myocardial Perfusion Imaging}-- Nuclear medicine imaging of heart muscle blood flow. Stress-rest protocol. Detects ischemia and infarction.

\medterm{Cardiac Gated SPECT}-- SPECT imaging synchronized with ECG. Assess ventricular function, wall motion, viability.

\medterm{Positron Emission Mammography (PEM)}-- PET imaging of breast. Higher sensitivity than mammography for dense breasts. Limited availability.

\medterm{Breast MRI}-- MRI of breast for high-risk screening, implant evaluation, staging breast cancer. High sensitivity, lower specificity.

\medterm{Diffusion-Weighted Imaging (DWI)}-- MRI technique measuring water molecule diffusion. Restricted diffusion bright. Used for stroke, abscess, tumor characterization.

\medterm{Perfusion Imaging}-- Imaging tissue blood flow. CT perfusion, MR perfusion, PET perfusion. Used for stroke, tumor characterization.

\medterm{Functional MRI (fMRI)}-- MRI measuring brain activity by detecting blood flow changes. Used for pre-surgical planning, research.

\medterm{Magnetic Resonance Spectroscopy (MRS)}-- MRI technique measuring metabolic compounds. Used for tumor characterization, metabolic disorders.

\medterm{CT Perfusion}-- CT technique measuring brain or organ perfusion. Used for acute stroke, tumor evaluation.

\medterm{Dual-Energy CT}-- CT using two different energy spectra. Material differentiation, virtual non-contrast images, gout detection.

\medterm{Cardiac CT}-- CT imaging of heart and coronary arteries. Coronary calcium scoring, coronary CT angiography, cardiac function assessment.

\medterm{CT Coronary Angiography}-- CT imaging of coronary arteries after contrast. Detects stenosis, plaque. Alternative to invasive angiography in low-intermediate risk.

\medterm{Coronary Calcium Scoring}-- Non-contrast CT measuring calcified plaque in coronary arteries. Agatston score predicts cardiovascular risk.

\medterm{Transcranial Doppler}-- Ultrasound measuring blood flow velocity in cerebral arteries. Used for vasospasm detection, stroke risk assessment.

\medterm{Carotid Duplex}-- Ultrasound combining B-mode and Doppler imaging of carotid arteries. Assesses stenosis, plaque characterization.

\medterm{Renal Artery Doppler}-- Ultrasound assessing renal artery stenosis. Screening for renovascular hypertension.

\medterm{Lower Extremity Venous Doppler}-- Ultrasound assessing deep vein thrombosis. Compression ultrasound, Doppler flow assessment.

\medterm{Upper Extremity Venous Doppler}-- Ultrasound assessing upper extremity DVT. Common in catheter-related thrombosis.

\medterm{Abdominal Aortic Ultrasound}-- Screening for abdominal aortic aneurysm. One-time screening in men 65-75 who smoked.

\medterm{Obstetric Ultrasound}-- Ultrasound imaging during pregnancy. First trimester: dating, viability. Second trimester: anatomy survey. Third trimester: growth, position.

\medterm{Fetal Echocardiography}-- Ultrasound imaging of fetal heart. Evaluates congenital heart disease.

\medterm{Doppler Velocimetry}-- Ultrasound measuring blood flow velocity in umbilical and fetal vessels. Assesses fetal well-being.

\medterm{Elastography}-- Imaging tissue stiffness. Used for liver fibrosis assessment, breast lesion characterization, thyroid nodule evaluation.

\medterm{Contrast-Enhanced Ultrasound}-- Ultrasound using microbubble contrast agents. Improves lesion characterization, perfusion assessment.

\medterm{Image-Guided Biopsy}-- Biopsy performed under imaging guidance. CT, ultrasound, MRI, or fluoroscopy. Improves accuracy, safety.

\medterm{Ablation Therapy}-- Image-guided destruction of tissue. Radiofrequency ablation, cryoablation, microwave ablation. Used for tumors, arrhythmias, pain.

\medterm{Embolization Therapy}-- Image-guided occlusion of vessels. Used for hemorrhage control, tumor devascularization, aneurysm treatment, varices.

\medterm{Thrombectomy}-- Image-guided removal of thrombus. Mechanical, pharmacological, or combined. Used for stroke, DVT, PE.

\medterm{Stent Placement}-- Image-guided deployment of stent. Vascular, biliary, ureteral, esophageal. Maintains patency.

\medterm{Angioplasty}-- Image-guided balloon dilation of stenosis. Coronary, peripheral, renal, carotid. Often combined with stenting.

\medterm{Drainage Procedures}-- Image-guided drainage of fluid collections. Abscess, hematoma, biloma, urinoma. CT or ultrasound guidance.

\medterm{Vertebroplasty}-- Percutaneous injection of bone cement into fractured vertebra. Provides pain relief, stabilization.

\medterm{Kyphoplasty}-- Similar to vertebroplasty but uses balloon to create cavity before cement injection.

\medterm{Discography}-- Contrast injection into intervertebral disc. Evaluates disc pathology. Controversial, limited use.

\medterm{Neurography}-- MRI technique visualizing peripheral nerves. Used for nerve tumors, entrapment, trauma.

\medterm{MR Neurography}-- See neurography.

\medterm{Lymphoscintigraphy}-- Nuclear medicine imaging of lymphatic system. Uses Tc-99m nanocolloid. Sentinel lymph node biopsy, lymphedema evaluation.

\medterm{Lymphangiography}-- See lymphangiography.

\medterm{Sialography}-- See sialography.

\medterm{Pneumography}-- Contrast imaging of body cavities. Pneumoencephalography obsolete. Pneumoperitoneum for laparoscopy.

\medterm{Arthrography}-- See arthrography.

\medterm{Synovography}-- Contrast imaging of joint cavity. Evaluates synovial pathology.

\medterm{Dacryocystography}-- Contrast imaging of lacrimal drainage system. Evaluates obstruction.

\medterm{Tubography}-- Contrast imaging of fallopian tubes. Evaluates patency.

\medterm{Hysterosalpingo-Contrast-Sonography (HyCoSy)}-- Ultrasound evaluation of tubal patency using contrast. Alternative to hysterosalpingography.

\medterm{Saline Infusion Sonohysterography}-- Ultrasound evaluation of uterus using saline contrast. Evaluates intrauterine pathology.

\medterm{Hysterosalpingography}-- See hysterosalpingography.

\medterm{Chelation Study}-- Nuclear medicine study using chelating agents. Evaluates renal function, obstruction.

\medterm{Mebrofenin Scan}-- Hepatobiliary imaging using Tc-99m mebrofenin. Evaluates liver function, biliary obstruction.

\medterm{Disofenin Scan}-- Hepatobiliary imaging using Tc-99m disofenin. Similar to mebrofenin scan.

\medterm{Iminoacetic Acid Scan}-- Renal imaging using Tc-99m iminodiacetic acid. Evaluates renal function, secretion.

\medterm{Sulfur Colloid Scan}-- Nuclear medicine imaging using Tc-99m sulfur colloid. Evaluates liver, spleen, bone marrow, lymph nodes.

\medterm{Red Blood Cell Labeling}-- Nuclear medicine technique labeling RBCs with Tc-99m. Used for GI bleeding detection, cardiac blood pool imaging.

\medterm{White Blood Cell Labeling}-- Nuclear medicine technique labeling WBCs with Tc-99m. Used for infection localization.

\medterm{Gallium Scan}-- Nuclear medicine imaging using Ga-67. Detects inflammation, infection, tumors. Largely replaced by FDG-PET.

\medterm{Indium-111 Scan}-- Nuclear medicine imaging using In-111. Used for WBC labeling, octreotide scanning.

\medterm{Octreoscan}-- Nuclear medicine imaging using In-111 octreotide. Detects neuroendocrine tumors.

\medterm{MIBG Scan}-- Nuclear medicine imaging using I-123 or In-111 MIBG. Detects pheochromocytoma, neuroblastoma.

\medterm{Thallium Scan}-- Nuclear medicine imaging using Tl-201. Myocardial perfusion, tumor detection. Largely replaced by Tc-99m agents.

\medterm{Technetium-99m}-- See nuclear medicine.

\medterm{Fludeoxyglucose F-18}-- See nuclear medicine.

\medterm{Radiopharmaceutical}-- See nuclear medicine.

\medterm{PET/CT}-- See nuclear medicine.

\medterm{PET/MRI}-- See nuclear medicine.

\medterm{MRI Contrast}-- See contrast agent.

\medterm{CT Contrast}-- See contrast agent.

\medterm{Ultrasound Contrast}-- See contrast agent.

\medterm{Barium Study}-- X-ray imaging using barium sulfate contrast. Barium swallow, barium enema, upper GI series. Evaluates GI tract morphology and function.

\medterm{Gastrografin Study}-- X-ray imaging using water-soluble contrast. Used when barium contraindicated (perforation risk). Less detailed than barium.

\medterm{Upper GI Series}-- X-ray imaging of esophagus, stomach, duodenum using barium. Evaluates motility, ulcers, masses, strictures.

\medterm{Small Bowel Series}-- X-ray imaging of small intestine using barium. Evaluates Crohn disease, tumors, obstructions.

\medterm{Small Bowel Follow-Through}-- Extended upper GI series tracking barium through small bowel.

\medterm{Barium Enema}-- X-ray imaging of colon using barium enema. Evaluates polyps, tumors, diverticula, inflammatory bowel disease.

\medterm{Defecography}-- X-ray imaging of defecation using barium. Evaluates pelvic floor dysfunction, rectocele, intussusception.

\medterm{Swallow Study}-- X-ray imaging of swallowing using barium. Evaluates dysphagia, aspiration, structural abnormalities.

\medterm{Videofluoroscopic Swallow Study}-- Real-time X-ray imaging of swallowing. Detailed assessment of oropharyngeal phase.

\medterm{Esophageal Manometry}-- Measurement of esophageal pressures. Evaluates motility disorders.

\medterm{pH Monitoring}-- Measurement of esophageal pH. Evaluates gastroesophageal reflux disease.

\medterm{Anorectal Manometry}-- Measurement of anal pressures. Evaluates fecal incontinence, constipation.

\medterm{Electroencephalography (EEG)}-- Recording of electrical activity from scalp electrodes. Used for seizure disorders, encephalopathy, brain death.

\medterm{Electromyography (EMG)}-- Recording of electrical activity from muscles. Used for neuropathy, myopathy, radiculopathy.

\medterm{Nerve Conduction Studies}-- Measurement of nerve conduction velocity and amplitude. Used for peripheral neuropathy, entrapment neuropathies.

\medterm{Evoked Potentials}-- Measurement of electrical responses to stimuli. Visual, auditory, somatosensory. Used for demyelinating diseases.

\medterm{Polysomnography}-- Sleep study measuring brain activity, breathing, oxygen saturation, muscle activity. Diagnoses sleep apnea, parasomnias.

\medterm{Holter Monitoring}-- Continuous ECG recording for 24-48 hours. Detects arrhythmias, ischemia.

\medterm{Event Monitor}-- Extended ECG monitoring for infrequent symptoms. Worn for weeks to months.

\medterm{Implantable Loop Recorder}-- Subcutaneous ECG monitor for years. Detects infrequent arrhythmias.

\medterm{Exercise Stress Test}-- ECG monitoring during exercise. Detects coronary artery disease.

\medterm{Pharmacological Stress Test}-- Stress testing using medications (dobutamine, adenosine, regadenoson). For patients unable to exercise.

\medterm{Stress Echocardiography}-- See stress echocardiography.

\medterm{Nuclear Stress Test}-- Myocardial perfusion imaging during stress. Detects coronary artery disease.

\medterm{Coronary Angiography}-- Invasive X-ray imaging of coronary arteries. Gold standard for diagnosing coronary artery disease.

\medterm{Venography}-- Contrast imaging of veins. Evaluates deep vein thrombosis, venous obstruction.

\medterm{Cerebral Angiography}-- Contrast imaging of cerebral vessels. Evaluates aneurysms, AVMs, stenosis. Invasive procedure.

\medterm{Spinal Angiography}-- Contrast imaging of spinal vessels. Evaluates vascular malformations.

\medterm{Pulmonary Angiography}-- Contrast imaging of pulmonary arteries. Gold standard for pulmonary embolism diagnosis. Largely replaced by CTPA.

\medterm{Aortography}-- Contrast imaging of aorta. Evaluates aneurysm, dissection, occlusion. Largely replaced by CTA/MRA.

\medterm{Renography}-- Nuclear medicine imaging of kidneys. Evaluates function, obstruction, blood flow.

\medterm{Renovascular Hypertension}-- Hypertension due to renal artery stenosis. Diagnosis: Doppler ultrasound, CTA, MRA, renography. Treatment: angioplasty, stenting, surgery.

\medterm{Hydronephrosis}-- Dilation of renal pelvis and calyces due to obstruction. Diagnosis: ultrasound, CT, MRI. Causes: stones, stricture, tumor, pregnancy.

\medterm{Renal Mass}-- Abnormality in kidney detected on imaging. Solid masses may be renal cell carcinoma. Cystic masses classified by Bosniak system.

\medterm{Bosniak Classification}-- System classifying renal cysts by CT features. Categories I-II benign, IIF follow-up, III-IV likely malignant.

\medterm{Liver Mass}-- Abnormality in liver detected on imaging. Benign: hemangioma, focal nodular hyperplasia, adenoma. Malignant: hepatocellular carcinoma, metastasis, cholangiocarcinoma.

\medterm{Pancreatic Mass}-- Abnormality in pancreas detected on imaging. Pancreatic adenocarcinoma most common malignancy. Neuroendocrine tumors also important.

\medterm{Splenomegaly}-- Enlarged spleen detected on imaging. Causes: portal hypertension, hematologic disorders, infections, infiltrative diseases.

\medterm{Hepatomegaly}-- Enlarged liver detected on imaging. Causes: fatty liver, hepatitis, congestion, infiltration, malignancy.

\medterm{Ascites}-- Fluid in peritoneal cavity detected on imaging. Causes: cirrhosis, malignancy, heart failure, nephrotic syndrome. Ultrasound sensitive.

\medterm{Pleural Effusion}-- Fluid in pleural space detected on imaging. Chest X-ray, ultrasound, CT. Transudative vs. exudative.

\medterm{Pneumothorax}-- Air in pleural space detected on imaging. Chest X-ray, ultrasound, CT. Tension pneumothorax emergency.

\medterm{Consolidation}-- Lung parenchyma filled with fluid, detected on imaging. Pneumonia, edema, hemorrhage, malignancy.

\medterm{Nodule}-- Small lung mass detected on imaging. Solitary pulmonary nodule evaluation: size, morphology, growth rate, PET avidity.

\medterm{Mass}-- Large lung lesion detected on imaging. Differentiate nodule (small) from mass (large). Malignancy more likely with mass.

\medterm{Infiltrate}-- Patchy lung opacity detected on imaging. Non-specific, may represent infection, edema, hemorrhage, inflammation.

\medterm{Ground-Glass Opacity}-- Hazy lung opacity on CT without obscuring vessels. Non-specific, may represent inflammation, edema, fibrosis, malignancy.

\medterm{Consolidation}-- See consolidation.

\medterm{Atelectasis}-- Collapse of lung tissue detected on imaging. Plate-like, lobar, or total lung collapse. Causes: obstruction, compression, scarring.

\medterm{Emphysema}-- Destruction of lung parenchyma detected on imaging. Centriacinar, panacinar. CT sensitive for detection.

\medterm{Bronchiectasis}-- Permanent dilation of bronchi detected on imaging. CT diagnosis. Causes: infection, obstruction, immunodeficiency.

\medterm{Interstitial Lung Disease}-- Group of diseases affecting lung interstitium. CT patterns: usual interstitial pneumonia, nonspecific interstitial pneumonia, cryptogenic organizing pneumonia.

\medterm{Pulmonary Embolism}-- See emergency medicine.

\medterm{Aortic Dissection}-- See emergency medicine.

\medterm{Abdominal Aortic Aneurysm}-- See vascular surgery.

\medterm{Mesenteric Ischemia}-- Reduced blood flow to intestines. Acute: embolism, thrombosis. Chronic: atherosclerosis. CT angiography diagnosis.

\medterm{Splenic Rupture}-- Traumatic or spontaneous splenic injury. CT diagnosis. May require splenectomy.

\medterm{Renal Trauma}-- Kidney injury from trauma. CT staging guides management. Grades I-V based on injury severity.

\medterm{Hepatic Trauma}-- Liver injury from trauma. CT staging guides management. Grades I-V based on injury severity.

\medterm{Pancreatic Trauma}-- Pancreatic injury from trauma. CT diagnosis. May require surgery.

\medterm{Bowel Obstruction}-- Blockage of intestinal lumen. X-ray, CT diagnosis. Differentiate small vs. large bowel obstruction. Mechanical vs. paralytic ileus.

\medterm{Appendicitis}-- Inflammation of appendix. CT diagnosis. May present atypically. Surgical emergency.

\medterm{Diverticulitis}-- Inflammation of colonic diverticula. CT diagnosis. Complications: abscess, perforation, fistula.

\medterm{Cholecystitis}-- Inflammation of gallbladder. Ultrasound diagnosis. Stone, wall thickening, pericholecystic fluid.

\medterm{Cholangitis}-- Infection of bile ducts. CT, MRCP diagnosis. Charcot triad: fever, jaundice, right upper quadrant pain.

\medterm{Pancreatitis}-- Inflammation of pancreas. CT diagnosis. Edematous vs. necrotizing.

\medterm{Renal Colic}-- Pain from urinary stone. CT diagnosis. Stone location, size, obstruction assessment.

\medterm{Hydronephrosis}-- See hydronephrosis.

\medterm{Urolithiasis}-- See urology.

\medterm{Pyelonephritis}-- Kidney infection. CT may show striated nephrogram, perinephric stranding.

\medterm{Renal Abscess}-- Collection of pus in kidney. CT diagnosis. May require drainage.

\medterm{Xanthogranulomatous Pyelonephritis}-- See urology.

\medterm{Renal Cell Carcinoma}-- See urology.

\medterm{Bladder Cancer}-- See urology.

\medterm{Prostate Cancer}-- See oncology.

\medterm{Breast Cancer}-- See oncology.

\medterm{Lung Cancer}-- See oncology.

\medterm{Liver Cancer}-- See oncology.

\medterm{Pancreatic Cancer}-- See oncology.

\medterm{Colorectal Cancer}-- See oncology.

\medterm{Gastrointestinal Stromal Tumor (GIST)}-- Mesenchymal tumor of GI tract. CT diagnosis. KIT mutation common.

\medterm{Neuroendocrine Tumor}-- Tumor from neuroendocrine cells. CT, PET diagnosis. Octreotide scan useful.

\medterm{Lymphoma}-- See oncology.

\medterm{Leukemia}-- See oncology.

\medterm{Metastasis}-- See oncology.

\medterm{Tumor Staging}-- See oncology.

\medterm{Response Assessment}-- Evaluation of treatment response. RECIST criteria for solid tumors. PET response criteria for lymphoma.

\medterm{Surveillance Imaging}-- Repeat imaging to detect recurrence. Protocol varies by cancer type, stage, treatment.

\medterm{Incidentaloma}-- Incidental imaging finding of uncertain significance. May require further evaluation.

\medterm{Incidental Finding}-- See incidentaloma.

\medterm{Radiation Dose}-- See radiation safety.

\medterm{Effective Dose}-- Dose measure accounting for tissue sensitivity. Measured in millisieverts (mSv).

\medterm{Contrast Reaction}-- Adverse reaction to contrast agent. Mild: nausea, urticaria. Severe: anaphylaxis, bronchospasm. Pre-medication for prior reactions.

\medterm{Nephrogenic Systemic Fibrosis}-- See gadolinium-induced nephrogenic systemic fibrosis.

\medterm{Extravasation}-- Contrast leakage into surrounding tissue during injection. Pain, swelling, potential tissue necrosis.

\medterm{Image Quality}-- Assessment of diagnostic adequacy. Factors: resolution, contrast, noise, artifacts.

\medterm{Resolution}-- Ability to distinguish small structures. Spatial, temporal, contrast resolution.

\medterm{Signal-to-Noise Ratio}-- Measure of image quality. Higher SNR = better image.

\medterm{Artifacts}-- See artifacts.

\medterm{Motion Artifact}-- Image degradation from patient movement. Reduced by faster acquisition, breath-hold, sedation.

\medterm{Metal Artifact}-- Image degradation from metal implants. Reduced by metal artifact reduction sequences (MRI), iterative reconstruction (CT).

\medterm{Beam Hardening Artifact}-- CT artifact from dense structures. Reduced by beam hardening correction algorithms.

\medterm{Partial Volume Artifact}-- CT artifact from voxel containing multiple tissues. Reduced by thinner slices.

\medterm{Aliasing Artifact}-- Ultrasound artifact from high velocities. Reduced by increasing scale, using lower frequency transducer.

\medterm{Reverberation Artifact}-- Ultrasound artifact from multiple reflections. Reduced by changing angle, using tissue harmonic imaging.

\medterm{Shadowing Artifact}-- Ultrasound artifact from attenuating structures. Enhanced by harmonic imaging.

\medterm{Enhancement Artifact}-- Ultrasound artifact from fluid-filled structures. Enhanced by harmonic imaging.

\medterm{Compression Ultrasound}-- Ultrasound technique applying pressure. Used for DVT diagnosis, testicular evaluation.

\medterm{Elastography}-- See elastography.

\medterm{Three-Dimensional Ultrasound}-- Ultrasound creating 3D images. Used in obstetrics, gynecology.

\medterm{Four-Dimensional Ultrasound}-- Real-time 3D ultrasound. Used in obstetrics for fetal imaging.

\medterm{Image-Guided Therapy}-- See interventional radiology.

\medterm{Percutaneous Drainage}-- See drainage procedures.

\medterm{Percutaneous Biopsy}-- See image-guided biopsy.

\medterm{Stereotactic Biopsy}-- Biopsy using stereotactic guidance. Breast, brain, lung.

\medterm{Core Needle Biopsy}-- Biopsy using core needle. CT, ultrasound, or MRI guidance.

\medterm{Fine Needle Aspiration}-- Biopsy using fine needle. Cytology diagnosis. Ultrasound or CT guidance.

\medterm{Excisional Biopsy}-- Surgical removal of entire lesion. Diagnosis and treatment.

\medterm{Incisional Biopsy}-- Surgical removal of part of lesion. Diagnosis before definitive treatment.

\medterm{Image-Guided Ablation}-- See ablation therapy.

\medterm{Radiofrequency Ablation}-- See ablation therapy.

\medterm{Cryoablation}-- See ablation therapy.

\medterm{Microwave Ablation}-- See ablation therapy.

\medterm{High-Intensity Focused Ultrasound (HIFU)}-- Ultrasound ablation using focused energy. Used for uterine fibroids, prostate cancer, bone metastases.

\medterm{Brachytherapy}-- Internal radiation therapy. Sources placed inside or near tumor. Used for prostate, cervical, breast cancer.

\medterm{External Beam Radiation Therapy}-- Radiation from outside body. 3D conformal, IMRT, stereotactic body radiation therapy (SBRT).

\medterm{Intensity-Modulated Radiation Therapy (IMRT)}-- Radiation therapy modulating beam intensity. Improves conformity, spares normal tissue.

\medterm{Stereotactic Radiosurgery}-- Single high-dose radiation to small target. Used for brain metastases, AVMs, acoustic neuroma.

\medterm{Proton Therapy}-- Radiation using protons instead of photons. Bragg peak allows precise dose delivery. Used for pediatric tumors, skull base tumors.

\medterm{Carbon Ion Therapy}-- Radiation using carbon ions. Higher biological effectiveness. Limited availability.

\medterm{Radiation Oncology}-- See radiation therapy.

\medterm{Radiation Physics}-- Physics of radiation imaging and therapy.

\medterm{Radiation Biology}-- Biological effects of radiation. Deterministic vs. stochastic effects.

\medterm{Radiation Protection}-- See radiation safety.

\medterm{Radiation Safety Officer}-- Person responsible for radiation safety program.

\medterm{Radiation Monitoring}-- Measurement of radiation exposure. Personal dosimeters, area monitors.

\medterm{Radiation Emergency}-- See disaster preparedness.

\medterm{Nuclear Medicine}-- See nuclear medicine.

\medterm{Diagnostic Radiology}-- See radiology.

\medterm{Interventional Radiology}-- See interventional radiology.

\medterm{Radiation Oncology}-- See radiation oncology.

\medterm{Medical Physics}-- See radiation physics.

\medterm{Radiologic Technologist}-- Professional performing imaging examinations.

\medterm{Radiologist}-- Physician interpreting imaging studies.

\medterm{Interventional Radiologist}-- Physician performing image-guided procedures.

\medterm{Radiation Oncologist}-- Physician providing radiation therapy.

\medterm{Medical Physicist}-- Expert in radiation physics and safety.

\medterm{ PACS}-- Picture Archiving and Communication System. Digital image storage and retrieval.

\medterm{DICOM}-- Digital Imaging and Communications in Medicine. Standard for medical image exchange.

\medterm{RIS}-- Radiology Information System. Workflow management for radiology department.

\medterm{EMR}-- Electronic Medical Record. Patient medical records.

\medterm{EHR}-- Electronic Health Record. Broader than EMR, includes population health.

\medterm{Tele radiology}-- Remote interpretation of imaging studies.

\medterm{Second Opinion}-- Review by another radiologist. Common for difficult cases, cancer staging.

\medterm{Quality Assurance}-- See quality improvement.

\medterm{Quality Control}-- See quality improvement.

\medterm{Accreditation}-- Formal recognition of compliance with standards. ACR accreditation for facilities.

\medterm{Board Certification}-- Formal recognition of radiology expertise. American Board of Radiology.

\medterm{Continuing Medical Education}-- See continuing education.

\medterm{Research}-- See research.

\medterm{Clinical Trials}-- See clinical trial.

\medterm{Evidence-Based Medicine}-- See evidence-based medicine.

\medterm{Guidelines}-- See clinical guidelines.

\medterm{Consensus}-- See guidelines.

\medterm{Standard of Care}-- See guidelines.

\medterm{Informed Consent}-- See informed consent.

\medterm{Patient Communication}-- See communication.

\medterm{Critical Results}-- See critical findings.

\medterm{Critical Findings}-- Findings requiring immediate communication. Examples: pneumothorax, brain herniation, aortic dissection.

\medterm{Incidental Findings}-- See incidental finding.

\medterm{Follow-Up Recommendations}-- Recommendations for further imaging or evaluation. Based on guidelines ( Fleischner Society for lung nodules).

\medterm{Reporting Standards}-- Standards for structured reporting. ACR guidelines, BI-RADS, LI-RADS, PI-RADS.

\medterm{Structured Reporting}-- Reporting using standardized templates. Improves consistency, completeness.

\medterm{Impression}-- Summary of key findings and recommendations in radiology report.

\medterm{Findings}-- Detailed description of imaging observations in radiology report.

\medterm{Technique}-- Description of imaging protocol in radiology report.

\medterm{Indication}-- Clinical reason for imaging in radiology report.

\medterm{Comparison}-- Reference to prior imaging in radiology report.

\medterm{Conclusion}-- See impression.

\medterm{Summary}-- See impression.
